\practicechapter{第 3 章实践：数据表示、内存、指针和对象布局}

本章对应第一册第 3 章“数据表示、内存、指针和对象布局”。正文讲字节、类型、对象表示、对齐、指针和布局。本章把这些概念落到 ELF 解析器的核心实现：所有字段最终都来自一段字节，解析器必须按小端规则组合整数，并且每次读取前都要证明范围存在。

\practicesection{一、对应原书问题：字段不是自动长出来的}

ELF header 里的 \code{machine} 是一个 16 位整数，\code{entry} 是一个 64 位整数，section header 里还有 offset、size、flags 等字段。文件中没有“整数对象”，只有连续字节。解析器把这些字节解释成整数，靠的是格式合同和字节序规则。

源码里的 \code{read\_u16\_le}、\code{read\_u32\_le}、\code{read\_u64\_le} 就是本章主角。它们不依赖把字节强转成结构体指针，而是逐字节组合。这种写法繁琐一点，但能清楚处理未对齐、对象生命周期、别名和越界问题，适合教学和可移植解析器。

\practicesection{二、从特例推到规律：小端读取和范围检查}

先看特例：字节序列 \code{78 56 34 12} 按 little-endian 解释为 \code{0x12345678}，按 big-endian 则是 \code{0x78563412}。如果不先检查 ELF ident 的 endianness，就无法知道同一段字节该怎样解释。

再看范围特例：输入长度为 10，想从 offset 8 读取 8 字节。offset 本身没有越界，但 offset 加 size 已经超过输入末尾。源码里的 \code{has\_range} 用 \code{offset <= byte\_count && size <= byte\_count - offset} 避免加法溢出和越界读取。这比直接判断 \code{offset + size <= byte\_count} 更稳。

由这两个特例推出本章规律：底层解析先证明“能读”，再证明“该怎样读”，最后才讨论“读出来代表什么”。顺序一乱，就会把未定义行为、越界和格式错误混成一个难排查的问题。

\practicesection{三、实践任务：给读取函数补证据}

阅读 \filepath{src/executable\_evidence.cpp} 中的四个函数：\code{has\_range}、\code{read\_u16\_le}、\code{read\_u32\_le}、\code{read\_u64\_le}。然后在测试报告中写出三个样例：

\begin{enumerate}
  \item \code{01 00} 读成 \code{1}；
  \item \code{78 56 34 12} 读成 \code{0x12345678}；
  \item offset 接近输入末尾时必须拒绝读取。
\end{enumerate}

本卷当前没有把这些私有函数直接暴露给测试，这是一个工程选择：对外合同测试通过 synthetic ELF 间接覆盖它们。作为本章加分任务，可以讨论是否应该把小端读取抽到内部可测试模块。若抽出，必须避免扩大公开 API；测试可以放在同一个 translation unit 或通过小型内部 helper 完成。

\practicesection{四、验收：内存解释要可防错}

本章交付物是 \filepath{reports/ch03-bytes-layout.md}。通过标准如下：

\begin{enumerate}
  \item 能用自己的话解释 little-endian，不只写“小端就是反过来”。
  \item 能说明为什么本项目不把 ELF header 直接 \code{reinterpret\_cast} 成 C++ struct。
  \item 能解释 \code{has\_range} 为什么不用简单的 \code{offset + size} 判断。
  \item 能把 section 的 offset 和 size 看成“文件字节范围”，而不是抽象名词。
  \item 能指出 padding、alignment、对象生命周期这些 C++ 规则和文件格式解析之间的边界。
\end{enumerate}

这章学完后，读者再看到 section、symbol、entry 时，应该能自然追问：这个字段从哪些字节来，按什么规则解释，读之前怎样证明安全。
